\section{Finite Galois inertia at the repaired tame prime}
\label{sec:cubic-hsieh-tame-inertia}

\begin{corollary}[AC-G70-inertia: the first tame type is killed by an explicit finite CM cover]
\label{AC-G70-inertia}
Choose an absolutely irreducible classical point of AC-G70
in its sufficiently regular separable cone.  Let $\rho_x$
be its three-dimensional semisimple Galois representation,
with the conjugate-self-dual normalization of Hsieh \S8.3.
At the first added split prime $q$, its inertia action is
finite, factors through the cyclic group of order $p^r$
cut out by $\vartheta_q$, and has zero monodromy.

Let $F_\vartheta/\mathbb Q$ be the cyclic field cut out by
the rational Dirichlet character $\vartheta$, and set
$L=KF_\vartheta$.  Then $L$ is CM, $L/K$ is cyclic of
degree $p^r$, and $\rho_x|_{G_L}$ is absolutely irreducible
and unramified at all primes above $q$.
No unramifiedness or zero-monodromy claim at $q'$ or $269$
is included.  This finite cover preserves the Hodge--Tate
weights and does not resolve their ordered symmetric
assignment or the actual Hecke-support gate.
\end{corollary}
\begin{proof}
\proofstep{Identify the exact stable packet to which the primary theorem applies.}
We use Bella\"iche--Chenevier,
\emph{Formes non temp\'er\'ees pour $U(3)$ et conjectures
de Bloch--Kato}, Ann. Sci. ENS (4) 37 (2004),
\S3.2.1--\S3.2.3, printed pp.~617--619.
Their five-case classification assigns reducible
three-dimensional Galois representations to every
endoscopic case and to the stable non-tempered
one-dimensional packet.  Consequently the absolutely
irreducible Galois representation at our point belongs
to the stable tempered case.
The stable packet has cuspidal base change to
$\mathrm{GL}_3/K$: this is the precise alternative from
Rogawski Theorem~13.3.3(b) used in their Proposition~3.1
proof, printed p.~620; its only other case is the
one-dimensional packet, already excluded.
The cohomological stable packet transfer between
$U(2,1)$ and its form compact at infinity, with unchanged
finite local packets, is the published transfer used
in the same proof.  Thus their Proposition~3.2 applies
to this packet.  These packet and base-change results
are retained published inputs, not new classification
claims.

The normalization agrees for the needed inertia
assertion.  Bella\"iche--Chenevier \S3.2.2--\S3.2.3
explicitly take the conjugate-self-dual representation
$\rho_{p,\mathrm{Rog}}(1)$ in the stable case;
Hsieh \S8.3 takes $R_p(\pi_x)(1)$.
Equivalently, the two semisimple representations are
identified by their same almost-everywhere Frobenius
polynomials and Chebotarev.
A Tate twist would in any event be unramified at $q\ne p$.

\proofstep{Use the kernel assertion, not an unproved full parameter comparison.}
By AC-G67-type the automorphic local parameter at a
prime $w\mid q$ is a sum of three characters
$\phi_1\oplus\phi_2\oplus\phi_3$ whose restrictions
to local units are
\[
                 \vartheta_q^{-1},\quad1,\quad\vartheta_q.
                                                        \tag{G70.10}
\]
The normalization by a power of the absolute value does
not affect these unit characters.
Bella\"iche--Chenevier Proposition~3.2, printed
pp.~620--621, states precisely
\[
 \rho_x(I'_w)=1,
 \qquad
 I'_w=\operatorname{rec}_w
       \ker\bigl((\phi_1\oplus\phi_2\oplus\phi_3)
                                      |_{\mathcal O_{K_w}^{\times}}\bigr),
                                                        \tag{G70.11}
\]
with the corresponding preimage in inertia understood.
Its complete proof constructs auxiliary solvable totally
real base changes with the prescribed local splitting
field, preserves cuspidality at each step, and applies
Proposition~3.1 to the resulting unramified local
base change.  The stated kernel assertion is the input
used here; no equality of all ramified local parameters
is presumed.

In (G70.10) the common kernel is exactly
$\ker\vartheta_q$.  Therefore the actual inertia action
of $\rho_x$ factors through its cyclic quotient of order
$p^r$.  This also proves that its Weil--Deligne monodromy
is zero.  Indeed on an open subgroup of inertia the
local representation has the form $\exp(t_p(\cdot)N)$;
finite image becomes trivial on a smaller open subgroup,
where $t_p$ still has an infinite open image.  A nonzero
nilpotent $N$ would give a nonconstant polynomial
$\exp(tN)$, which cannot be the identity for infinitely
many $t$ in characteristic zero.  Hence $N=0$.
We have proved the actual kernel and finite inertia,
without needing to identify the three inertia
multiplicities individually.

\proofstep{Kill this inertia while preserving irreducibility.}
The rational character $\vartheta$ has odd order $p^r$,
so $\vartheta(-1)=1$ and $F_\vartheta$ is totally real.
Since $K$ is imaginary quadratic, their intersection
is $\mathbb Q$.  Hence $L=KF_\vartheta$ is CM and
$\operatorname{Gal}(L/K)$ is cyclic of order $p^r$.
At every prime above $q$ the restricted inertia is
inside $\ker\vartheta_q$, so (G70.11) makes
$\rho_x|_{G_L}$ unramified there.

For clarity, irreducibility here does not require an
extra residual hypothesis.  Restriction of a finite-dimensional
irreducible characteristic-zero representation to a
normal subgroup of finite index is semisimple: the
conjugates of any irreducible submodule span the whole
space, and a sum of irreducible submodules is semisimple.
By induction--restriction and the decomposition of the
regular representation of the finite cyclic quotient,
\[
 \operatorname{End}_{G_L}(\rho_x)
   \simeq\bigoplus_{\eta\in\widehat{\operatorname{Gal}(L/K)}}
                    \operatorname{Hom}_{G_K}(\rho_x,\rho_x\otimes\eta).
                                                        \tag{G70.12}
\]
A nonzero summand gives an isomorphism; determinants
then give $\eta^3=1$.  Since $\eta$ also has $p$-power
order and $p>3$, it is trivial.  Over the algebraically
closed coefficient field the endomorphism algebra in
(G70.12) thus has dimension one.  Semisimplicity makes
the restricted representation absolutely irreducible.
Restriction to a finite local extension leaves the
Hodge--Tate weights at the extended embedding unchanged,
so it supplies none of the missing ordered-weight bridge.
\end{proof}

\paragraph{Scope and geometric dictionary.}
This closes the finite-inertia requirement at the first
repaired prime for the chosen stable classical points,
and gives an explicit finite CM cover killing that
inertia.  It neither prescribes the new $q'$ type nor
removes possible repeated-character monodromy at $269$.
No size estimate for the enlarged arithmetic level,
no symmetric ordinary refinement, and no AC-G61 support
containment follows.  The cover and restriction are
ordinary finite-cover operations in the previously
audited SGA~1, Expos\'e~V, \S\S4 and~7 framework.
